\begin{table}[!htbp]\centering
\caption{The disappeared volume was real capacity, not low-intensity custodial care}
\label{tab:psych-intensity}
\small
\begin{tabular}{lr}
\toprule
Pre-closure mean length of stay (days) & \\
\midrule
\quad at the closing hospital & 24.0 \\
\quad at surviving hospitals & 23.1 \\
\midrule
Post-closure decline (exposed catchments) & \\
\quad inpatient psychiatric admissions & 63\% \\
\quad inpatient psychiatric bed-days & 66\% \\
\bottomrule
\multicolumn{2}{p{0.80\textwidth}}{\footnotesize Notes: Length of stay is admission-weighted, pooled over event years $-3$ to $-1$ in exposed catchments. Admissions at the closing hospitals were no shorter than at the hospitals that survived, and bed-days fall as steeply as admission counts---so the contraction removed full-intensity inpatient capacity, not residual custodial volume. Where that capacity went is only partly identifiable from inpatient records.}\\
\end{tabular}
\end{table}
